\section*{Resumen del curso}
\addcontentsline{toc}{section}{Resumen del curso}

Las matemáticas son el lenguaje en el que se escribe la física. En concreto, el cálculo y el álgebra tienen una capacidad expresiva tremenda, que los hacen idóneos para escribir las relaciones y dinámica de los observables físicos. Ya solo con esas áreas matemáticas se puede llegar muy lejos en física. Sin embargo, eventualmente es necesario rebuscar un poco más en el fondo de armario matemático. Una de las primeras ramas adicionales que cobra importancia es, sin duda, la geometría diferencial.

La geometría diferencial expande las técnicas del cálculo y el álgebra en tanto que permite aplicarlas en espacios con curvatura. En cursos iniciales de física solemos imaginar partículas que se mueven en un espacio euclídeo. Las trayectorias pueden ser curvas, claro, pero acontecen en un escenario plano. La aplicación más inmediata es la relatividad general, no en vano encontramos ya en la cultura popular el lema ``la gravedad es la curvatura del espacio-tiempo''. Sin embargo, la geometría diferencial juega un papel fundamental también en teorías como la electrodinámica o el propio modelo estándar de física de partículas; en teorías de campos gauge, en general. 

Por el teorema ``no free lunch'', las prestaciones superiores de la geometría diferencial no son gratis. Hay dificultades ligadas a los espacios con curvatura que requieren de un andamiaje conceptual más sofisticado. En este curso vamos a introducir, de la mano de la física, las ideas más importantes. No buscamos dar una visión global detallada, sino más bien pinceladas esenciales. Tampoco pretendemos mantener un rigor formal excesivo, aunque alguna sección pueda resultar un poco ``dura''. Pondremos el foco donde nos convenga, barriendo muchos detalles técnicos bajo la alfombra. De todos modos, a fin de paliar las carencias de estas notas, hemos incluido referencias didácticas para profundizar. Asimismo, las explicaciones van acompañadas de numerosos ejemplos y figuras. Algunas de las imágenes, las tridimensionales, son interactivas si se usa el lector de \texttt{pdf} de Adobe. 

El objetivo principal de estas notas es allanar el camino para el estudio de la relatividad general. No obstante, comenzar hablando de relatividad general es una receta para el desastre. En su lugar, conviene motivar y construir los fundamentos de la geometría diferencial desde teorías más controladas. Así comienza la \cref{sec.1motivacion}, que se apoya en la intuición para orquestar el primer encuentro con la geometría diferencial en el contexto de la mecánica clásica. Finalizamos la misma sección con un repaso de relatividad especial, enfatizando las ideas de geometría diferencial. 

A partir ahí, comienza la parte más formal del curso. El tratamiento en las siguientes secciones es más matemático-geométrico de lo que es habitual para un curso dirigido a físicos, donde normalmente se pondría un énfasis mucho mayor en la aplicación directa a la relatividad general. A cambio, estas notas pretenden ofrecer una imagen más clara y completa de los conceptos fundamentales, su significado geométrico y las relaciones entre ellos.

En la \cref{sec.2tensores} se introduce rigurosamente el concepto de tensor desde el álgebra. Comenzamos con un recordatorio de las bases del álgebra lineal, para después definir las formas lineales y multilineales. Con ese objetivo, abandonaremos explícitamente la notación de arrays que tan interiorizada está. La sustituiremos por la notación de índices, mucho más indicada para hablar de multilinealidad. Por supuesto, es imprescindible para entender cualquier texto de relatividad especial, pero la agilidad en el manejo de índices también es tremendamente útil en mecánica cuántica.

La \cref{sec.2tensores} es muy interesante por sí misma, pero la idea de este curso es aprender a hacer cálculo y álgebra en espacios con curvatura conocidos como variedades diferenciales. El espacio-tiempo de la relatividad es una de esas variedades diferenciales, por ejemplo, aunque también lo es una simple esfera. En la \cref{sec.1motivacion} ya se analizan algunos ejemplos con el apoyo de la intuición y la física, pero en la \cref{sec.3variedades} vamos más al detalle acerca de qué es una variedad diferencial, cómo manejarla y cómo definir tensores sobre ella. Por último, la \cref{sec.4curvatura} establece las bases de la descripción matemática de la curvatura. Para ello, definimos nociones como la conexión afín, la derivada covariante o los tensores de curvatura. 
